This activity covers user stories US-004, US-005 and US-041. Research projects and saved queries give every dataset its research context: a collection run always belongs to a project and starts from one of its saved queries.

\subsection{Research projects}

A research project has a name, a description, a default language, a status and an owner, and project names are unique per owner. The \emph{Projects} section lists the owner's projects, creates new ones and selects the active project, whose queries, runs and results the other sections show. The API allows projects to be listed, created, read and updated; it does not delete them, because deleting a project cascades to its saved queries, runs and request logs, which form the methodological record of its datasets. Deletion remains available to the administrator through the Django administration site.

Every endpoint below a project resolves the project among those of the current owner, and a project of another owner returns 404, the same answer as a project that does not exist, so the API reveals nothing about other users' data. The application has no login yet: requests without an authenticated session are attributed to a configured local owner. This is sufficient for a tool that listens only on the local machine, but authentication and rate limiting are required before any shared deployment, because every caller can spend the operator's LLM and YouTube quotas.

\subsection{Saved queries}

A saved query belongs to a project and stores the collection criteria: a name, the natural-language prompt when there was one, the search term, the publication date range, the region and language, the number of videos to discover, the optional numbers of comments per video and replies per comment, notes and the complete \texttt{search.\allowbreak list} parameters, together with the query source and, for generated queries, the LLM provenance described in Activity~\ref{act:2.2}. Query names are unique for the same creator within a project.

Saved queries can be created and listed but not edited through the API. A changed criterion is saved as a new query, so the relation between a run and the query that started it stays unambiguous, in addition to the parameter snapshot each run keeps (Activity~\ref{act:2.1}).

The \emph{Queries} section contains the natural-language draft panel and the saved-query form. Ten advanced \texttt{search.\allowbreak list} filters, added in US-041 at the request of the thesis advisor, are grouped in a collapsed part of the form: channel, live event type, location and radius, safe search, captions, category, definition, duration and license. They are set by hand, are not filled by the LLM and are left untouched when a draft is applied; a blank filter is not sent. They are stored inside the query's parameters, with no new database column, and are therefore copied into each run's snapshot like any other parameter. Each saved query is shown as a card with its criteria, its advanced filters and a button that starts a collection run.
